% ============================================================================
% GEM SECTION 2 --- Method.
%
% REWRITTEN 2026-08-27 FROM TEXT SUPPLIED VERBATIM.  Transcribed as given; no
% sentence reworded, shortened or resequenced.  Page fitting is the author's
% call, not this file's.
%
% The change from the previous version is density, not length-for-its-own-sake:
% the earlier compression sent the LEARNING OBJECTIVE and the factorization of
% R_theta to the appendix, which left Section 3.1's NLL decomposition without
% the construction it decomposes.  Eqs. (2)-(4) are restored to the body for
% exactly that reason.
%
% Load-bearing objects kept in the body, in order:
%   A(x)            what is executable
%   p_theta(a|x)    what is learned at mark level
%   R_theta(y|x)    why probability lives on molecular outcomes
%   L_ref           what the model actually learns
%   h_b             what "future-aware" means
%   P^{z,b}         how purpose changes the frozen process
%   Proposition 2   what the exact construction guarantees
%
% Still deliberately OUT of the body: SMC mechanics, Monte Carlo target
% estimators, architecture, full operator tables, ring-builder implementation,
% frontier algorithms, cache/inference machinery, proof details.
%
% APPENDIX POINTERS BELOW ARE AS SUPPLIED (N, O, L, G, Q).  They are transcribed
% verbatim and do NOT yet correspond to \label{}s in gem_appendix.tex; wire them
% before submission.
% ============================================================================

\section{COMPOSE: an executable molecular process}
\label{sec:method}

COMPOSE represents molecular design as control of a learned stochastic process
over complete molecules (Figure~\ref{fig:ppp}). The construction separates
three roles. An exact rewrite executor determines which molecular transitions
are \textbf{possible}. A goal-independent reference model learns which
executable transitions are \textbf{plausible}. Task-specific control then
reweights this frozen process according to which molecular futures are
\textbf{purposeful} for the current objective.

\subsection{Executable molecular state space}
\label{sec:gem-support}

Let
\[
  \mathcal X=\bigcup_{n=1}^{n_{\max}}\mathcal X_n
\]
denote the molecular state space, where $\mathcal X_n$ contains connected
molecular graphs with $n$ heavy atoms, declared atom and bond attributes, and
valid valences. Every $x\in\mathcal X$ is therefore a complete molecule rather
than a masked, corrupted, or partially specified state.

Transitions are typed graph rewrites. An edit mark $a$ specifies a rewrite
family, molecular operands, and a typed payload, and $T(x,a)$ denotes the
resulting graph. From state $x$, the executor exposes
\begin{equation}
  \mathcal A(x)
  =
  \left\{
    a:\; a\text{ matches } x
    \;\text{and}\;
    T(x,a)\text{ passes all executor guards}
  \right\}.
  \label{eq:gem-editset}
\end{equation}
The guards enforce the declared atom and bond vocabularies, valence
constraints, connectivity, and molecular-size bounds before an edit is
committed. Inadmissible states therefore have zero transition support rather
than being discouraged by a learned penalty.

The rewrite language contains atom insertion and deletion together with atom-,
bond-, and ring-level transformations (Appendix~N). Insertion and deletion
connect $\mathcal X_n$ to $\mathcal X_{n+1}$ and $\mathcal X_{n-1}$, while the
remaining transformations restructure a molecule at fixed size. COMPOSE is
therefore trans-dimensional: trajectories can grow, shrink, and reorganize
molecular graphs within one state space.

\begin{proposition}[Executable-support closure]
\label{prop:gem-closure}
For every $x\in\mathcal X$ and $a\in\mathcal A(x)$,
\[
  T(x,a)\in\mathcal X .
\]
\end{proposition}

Hence any stochastic process supported on the executable edit sets
$\{\mathcal A(x)\}_{x\in\mathcal X}$ remains in the declared molecular state
space. The proof and full rewrite semantics are given in Appendix~O.

The executor determines which molecular changes are possible. It does not
determine how probability should be distributed among them.

\subsection{Learning a molecular reference process}
\label{sec:gem-reference}

We learn a goal-independent stochastic law over executable edits. Writing a
mark as $a=(k,o,\eta)$, with rewrite family $k$, operands $o$, and typed payload
$\eta$, the reference model factorizes
\begin{equation}
  p_\theta(a\mid x)
  =
  p_\theta(k\mid x)\,
  p_\theta(o,\eta\mid k,x),
  \qquad a\in\mathcal A(x).
  \label{eq:gem-factorization}
\end{equation}
Normalization is restricted to candidates admitted by the executor. The model
receives no design objective, remaining budget, or downstream task information.
Thus $p_\theta$ models molecular transition structure independently of the
purpose for which the process will later be controlled.

The stochastic object used for control is a distribution over \textbf{molecular
successors}, not raw edit marks. Different marks can produce the same molecule
because of graph symmetries, equivalent matches, or alternative internal rewrite
descriptions. For a canonical successor $y$, define its executable fiber
\[
  \mathcal G_y(x)
  =
  \{a\in\mathcal A(x): T(x,a)\simeq y\},
\]
where $\simeq$ denotes registered molecular canonicalization. Let $\mathcal
S(x)$ denote the distinct non-self canonical successors of $x$. We define
\begin{equation}
  R_\theta(y\mid x)
  =
  \frac{
    \sum_{a\in\mathcal G_y(x)} p_\theta(a\mid x)
  }{
    \sum_{y'\in\mathcal S(x)}
    \sum_{a\in\mathcal G_{y'}(x)} p_\theta(a\mid x)
  },
  \qquad y\in\mathcal S(x).
  \label{eq:gem-rtheta}
\end{equation}
This quotient is important. A molecular transition does not become more probable
merely because symmetry gives the executor several syntactic ways to describe
it. State-level control therefore acts only after canonicalization.

We train $R_\theta$ from executable teacher transitions. If $y_i^\star$ is the
realized canonical successor of state $x_i$, the reference objective is
\begin{equation}
  \mathcal L_{\mathrm{ref}}(\theta)
  =
  -\frac{1}{N}\sum_{i=1}^{N}\log R_\theta(y_i^\star\mid x_i).
  \label{eq:gem-refloss}
\end{equation}
Every teacher transition is replayed through the same executor used at
inference. Targets and compiled programs are used to construct supervision but
are not supplied to the reference model at inference. We learn this reference
process once and keep it fixed across downstream objectives.

Equations~\eqref{eq:gem-factorization}--\eqref{eq:gem-refloss} also motivate the
decomposition in Section~\ref{sec:gem-exp-rtheta}: the family head learns
\textbf{which kind of edit} is appropriate, while the within-family law learns
\textbf{which molecular successor} is appropriate for the current state.

\subsection{Finite-horizon control of the frozen process}
\label{sec:gem-control}

Task information enters only after the molecular reference process has been
learned. Let $z$ denote a design objective and $g_z:\mathcal X\rightarrow
\mathbb R_{\geq0}$ its terminal desirability. For $b$ remaining molecular
transitions, define
\begin{equation}
  h_0(x,z)=g_z(x),
  \qquad
  h_b(x,z)
  =
  \sum_{y\in\mathcal S(x)} R_\theta(y\mid x)\, h_{b-1}(y,z).
  \label{eq:gem-backward}
\end{equation}
Equivalently,
\[
  h_b(x,z)
  =
  \mathbb E_{R_\theta}\!\left[ g_z(X_b)\mid X_0=x \right].
\]
This distinction is central. A property model scores \textbf{the molecule
currently in hand}. The backward value $h_b$ instead scores \textbf{what remains
reachable from that molecule within the remaining budget}.

For $h_b(x,z)>0$, the finite-horizon Doob transform defines
\begin{equation}
  P^{z,b}(y\mid x)
  =
  R_\theta(y\mid x)\,
  \frac{h_{b-1}(y,z)}{h_b(x,z)} .
  \label{eq:gem-central}
\end{equation}
The controller therefore changes probability only among molecular transitions
already supported by the frozen reference process.

\begin{proposition}[Finite-horizon molecular control]
\label{prop:gem-doob}
Starting from $x_0$ with budget $K$, applying
Equation~\eqref{eq:gem-central} at each step yields
\[
  \Pr\nolimits_{P^z}(X_K=y\mid X_0=x_0)
  =
  \Pr\nolimits_{R_\theta}(X_K=y\mid X_0=x_0)\,
  \frac{g_z(y)}{h_K(x_0,z)} .
\]
\end{proposition}

Thus exact backward values realize the $g_z$-tilted terminal law while
introducing no transition outside the executable reference support. The proof
and precise terminal-state assumptions are given in Appendix~O.

Exact backward computation requires enumerating the reachable molecular state
graph and is infeasible at molecular scale. We therefore learn a goal- and
budget-conditioned approximation
\begin{equation}
  h_\phi(x,z,b)\approx
  \mathbb E_{R_\theta}\left[ g_z(X_b)\mid X_0=x \right]
  \label{eq:gem-hphi}
\end{equation}
from finite-budget continuations of the frozen reference process. Replacing $h$
by $h_\phi$ in Equation~\eqref{eq:gem-central} gives the scalable future-value
controller. Approximation error may change which legal successor is preferred,
but it cannot introduce transitions outside the support of $R_\theta$. We
therefore reserve the exact terminal-reweighting guarantee for the exact
$h$-transform. Future-value supervision and sampling algorithms are given in
Appendices~G and~Q.

Because control is recomputed from the molecule actually reached, the same
process also supports intervention after execution begins. If an objective
changes from $z$ to $z'$ after reaching $x_\tau$, control simply resumes from
$x_\tau$ under $z'$; no parameter of $R_\theta$ changes.
Section~\ref{sec:gem-exp-lookahead} tests whether retaining this realized
molecular progress is useful in practice.

\paragraph{Executable transformations at longer time scales.}
\label{sec:gem-macros}
Some useful molecular changes require several primitive rewrites before their
value becomes visible. For molecular optimization, COMPOSE therefore exposes a
small set of temporally extended transformations, including linked and fused
ring construction and ring-local refinement. These operations introduce no
additional molecular support: they compile into the same primitive executable
rewrites, and every committed intermediate is checked by the same executor. The
controller chooses the transformation semantics, while the frozen reference law
ranks its legal molecular realizations. Full construction, refinement, and
qualification details are given in Appendix~L.
